\documentclass[11pt]{article}

\usepackage[margin=1in]{geometry}
\usepackage{amsmath,amssymb,mathtools}
\usepackage{booktabs}
\usepackage{array}
\usepackage{enumitem}
\usepackage{microtype}
\usepackage{xcolor}
\usepackage{fancyhdr}
\usepackage{lastpage}
\usepackage[hidelinks]{hyperref}

\setlength{\parindent}{0pt}
\setlength{\parskip}{5pt}
\setlist{nosep}
\renewcommand{\arraystretch}{1.25}
\newcommand{\entropy}{H}
\newcommand{\IG}{\operatorname{IG}}
\newcommand{\Rtwo}{R^2}
\newcommand{\E}{\mathbb{E}}

\pagestyle{fancy}
\fancyhf{}
\setlength{\headheight}{14pt}
\lhead{EECS 492 -- Homework 1}
\rhead{Written Solutions}
\cfoot{\thepage\ of \pageref{LastPage}}

\title{\vspace{-1.5cm}\textbf{EECS 492 -- Homework 1}\\[2pt]
\large Written Solutions}
\author{}
\date{Fall 2026}

\begin{document}
\maketitle
\vspace{-0.8cm}

\section{Designing Agents}

\subsection*{1(a) PEAS description}
The task environment contains a repeated sequence of tic-tac-toe games.  At the
end of a game, the board is cleared (or a fresh sheet is provided) and the next
game begins.  A PEAS description is therefore:

\begin{center}
\begin{tabular}{>{\raggedright\arraybackslash}p{0.17\textwidth}
                >{\raggedright\arraybackslash}p{0.75\textwidth}}
\toprule
\textbf{Component} & \textbf{Specification}\\
\midrule
Performance measure &
Correctly recognize the board geometry, the symbols $X$ and $O$, and empty
cells; make only legal moves in empty cells; choose a move that is optimal for
the agent's objective; correctly detect a win, loss, or tie; and reset and
repeat the procedure reliably over arbitrarily many games.  Secondary costs
include drawing errors, sensing errors, unnecessary pen motion, and time.\\
\addlinespace
Environment &
A physical sheet of paper containing a $3\times3$ grid, the current markings,
the pen and the board location, an opposing human or robot, and the sequence of
past game outcomes.  Between games the board is reset.\\
\addlinespace
Actuators &
A pen-moving mechanism, a pen-lowering/raising or contact mechanism, and a
marker-drawing action that selects one of the nine cells.  The agent may also
signal the result of a game and request or perform a reset.\\
\addlinespace
Sensors &
A camera or other visual sensor for locating the grid, classifying $X$ versus
$O$, and detecting empty cells; position/contact feedback for the pen; and a
state-change detector for recognizing the opponent's move and the terminal
board state.\\
\bottomrule
\end{tabular}
\end{center}

\textbf{Assumptions.} The paper remains within the sensing workspace, the grid
is not intentionally destroyed, and illumination and handwriting vary within
the sensor's operating range.  The opponent is allowed to be either a person
or another agent and may choose any legal move.

\subsection*{1(b) Environment characteristics}
The classification below is for the physical robot described in the question,
not for an idealized symbolic program.  The distinction matters because sensing
and pen motion introduce uncertainty even though the game rules themselves are
simple.

{\footnotesize
\renewcommand{\arraystretch}{1.10}
\begin{center}
\begin{tabular}{>{\raggedright\arraybackslash}p{0.24\textwidth}
                >{\raggedright\arraybackslash}p{0.18\textwidth}
                >{\raggedright\arraybackslash}p{0.50\textwidth}}
\toprule
\textbf{Dimension} & \textbf{Classification} & \textbf{Justification}\\
\midrule
Observability & Partially observable &
The true board state is not directly available as symbols.  The agent must infer
it from images, and glare, occlusion, ambiguous handwriting, or an imperfect
classifier can hide whether a cell is empty or which marker it contains.\\
\addlinespace
Number of agents & Multi-agent &
The opponent selects actions that change the board and therefore changes the
agent's future choices.\\
\addlinespace
Transition model & Nondeterministic &
The opponent's action is not controlled by the agent, and sensing or drawing can
fail.  If all perception and actuation were assumed perfect and the opponent's
move were supplied as an action, the symbolic game transition would instead be
deterministic.\\
\addlinespace
Temporal dependence & Sequential &
Each move changes the legal action set and can determine later wins, losses, or
ties.  The complete task is a sequence of such games, with a reset between
episodes.\\
\addlinespace
Change while deliberating & Dynamic &
The opponent or the physical setup may change while the robot is sensing or
planning.  If the opponent is required to wait and the sheet is fixed during
the robot's turn, the game is only semi-dynamic during that turn; the physical
specification here permits changes, so ``dynamic'' is the conservative label.\\
\addlinespace
State and actions & Discrete at the game level &
There are finitely many cells, marker values, legal moves, and terminal results.
Pen trajectories and camera measurements are continuous physical quantities,
but they are abstracted into discrete game states and actions.\\
\addlinespace
Model knowledge & Known at the rule level; uncertain at the physical level &
The agent can be programmed with the tic-tac-toe rules and winning condition.
It may not know the exact camera noise, pen-placement error, lighting, or
opponent policy before calibration and experience.\\
\bottomrule
\end{tabular}
\end{center}
}

\section{K-Nearest Neighbors}

The query point is
\[
q=(8,0,3).
\]
For Trial $m$, only the first $m$ coordinates are used, so
\[
d_m(q,x_i)=\left(\sum_{j=1}^{m}(q_j-x_{ij})^2\right)^{1/2}.
\]

\subsection*{2(a) Distances}
The individual calculations are:

\begin{align*}
\text{Trial 1: }d_1(q,x_1)&=|8-5|=3,\\
d_1(q,x_2)&=|8-7|=1,\\
d_1(q,x_3)&=|8-0|=8,\\
d_1(q,x_4)&=|8-2|=6,\\
d_1(q,x_5)&=|8-9|=1,\\
d_1(q,x_6)&=|8-3|=5.
\end{align*}

For Trial 2, the first two coordinates give
\begin{align*}
 d_2(q,x_1)&=\sqrt{(8-5)^2+(0-2)^2}=\sqrt{13}=3.606,\\
 d_2(q,x_2)&=\sqrt{(8-7)^2+(0-6)^2}=\sqrt{37}=6.083,\\
 d_2(q,x_3)&=\sqrt{(8-0)^2+(0-6)^2}=\sqrt{100}=10.000,\\
 d_2(q,x_4)&=\sqrt{(8-2)^2+(0-7)^2}=\sqrt{85}=9.220,\\
 d_2(q,x_5)&=\sqrt{(8-9)^2+(0-2)^2}=\sqrt{5}=2.236,\\
 d_2(q,x_6)&=\sqrt{(8-3)^2+(0-5)^2}=\sqrt{50}=7.071.
\end{align*}

For Trial 3, all three coordinates give
\begin{align*}
 d_3(q,x_1)&=\sqrt{(8-5)^2+(0-2)^2+(3-8)^2}=\sqrt{38}=6.164,\\
 d_3(q,x_2)&=\sqrt{(8-7)^2+(0-6)^2+(3-1)^2}=\sqrt{41}=6.403,\\
 d_3(q,x_3)&=\sqrt{(8-0)^2+(0-6)^2+(3-0)^2}=\sqrt{109}=10.440,\\
 d_3(q,x_4)&=\sqrt{(8-2)^2+(0-7)^2+(3-4)^2}=\sqrt{86}=9.274,\\
 d_3(q,x_5)&=\sqrt{(8-9)^2+(0-2)^2+(3-4)^2}=\sqrt{6}=2.449,\\
 d_3(q,x_6)&=\sqrt{(8-3)^2+(0-5)^2+(3-8)^2}=\sqrt{75}=8.660.
\end{align*}

Thus the requested table is:
\begin{center}
\begin{tabular}{c|ccc}
\toprule
Point & Trial 1 & Trial 2 & Trial 3\\
\midrule
$x_1$ & 3.000 & 3.606 & 6.164\\
$x_2$ & 1.000 & 6.083 & 6.403\\
$x_3$ & 8.000 & 10.000 & 10.440\\
$x_4$ & 6.000 & 9.220 & 9.274\\
$x_5$ & 1.000 & 2.236 & 2.449\\
$x_6$ & 5.000 & 7.071 & 8.660\\
\bottomrule
\end{tabular}
\end{center}

\subsection*{2(b) $k=3$ classification}
The training labels are $y_1=A$, $y_2=B$, $y_3=A$, $y_4=B$,
$y_5=A$, and $y_6=B$.

\begin{itemize}
    \item \textbf{Trial 1.} The sorted distances are $1$ for $x_2$, $1$ for
    $x_5$, and $3$ for $x_1$ (the first two are tied).  The three labels are
    $B,A,A$, so the majority is $A$.  The nearest-neighbor set is
    \(\{x_2,x_5,x_1\}\), up to the order of the tied points.
    \item \textbf{Trial 2.} The three smallest distances are for
    $x_5,x_1,x_2$, with labels $A,A,B$.  The prediction is \(\boxed{A}\).
    \item \textbf{Trial 3.} The three smallest distances are for
    $x_5,x_1,x_2$, with labels $A,A,B$.  The prediction is \(\boxed{A}\).
\end{itemize}
Therefore, all three trials classify the query as $A$.

\section{Decision Tree}

For a label set $S$, use
\[
\entropy(S)=-\sum_{c\in\{\mathrm{Yes},\mathrm{No}\}}
 p_c\log_2 p_c,
\qquad
\IG(S,A)=\entropy(S)-\sum_{v\in\operatorname{Values}(A)}
 \frac{|S_v|}{|S|}\entropy(S_v).
\]
The data contain $6$ ``Yes'' labels and $4$ ``No'' labels.  Therefore,
\begin{align*}
\entropy(Y)
&=-\frac{6}{10}\log_2\frac{6}{10}-\frac{4}{10}\log_2\frac{4}{10}\\
&=0.970950594\approx 0.971.
\end{align*}

\subsection*{3(a) Information gain at the root}

\paragraph{Outlook.}
The Sunny branch has $(1\ \mathrm{Yes},3\ \mathrm{No})$, the Overcast branch
has $(2\ \mathrm{Yes},0\ \mathrm{No})$, and the Rain branch has
$(3\ \mathrm{Yes},1\ \mathrm{No})$.  Hence
\begin{align*}
\entropy(\mathrm{Sunny})&=-\tfrac14\log_2\tfrac14-	frac34\log_2\tfrac34
=0.811278,\\
\entropy(\mathrm{Overcast})&=0,\\
\entropy(\mathrm{Rain})&=-\tfrac34\log_2\tfrac34-	frac14\log_2\tfrac14
=0.811278,\\
\IG(Y,\mathrm{Outlook})
&=0.970951-\tfrac4{10}(0.811278)-\tfrac2{10}(0)-\tfrac4{10}(0.811278)\\
&=\boxed{0.322}.
\end{align*}

\paragraph{Humidity.}
The High branch has $(2\ \mathrm{Yes},3\ \mathrm{No})$, while the Normal branch
has $(4\ \mathrm{Yes},1\ \mathrm{No})$.  Thus
\begin{align*}
\entropy(\mathrm{High})&=-\tfrac25\log_2\tfrac25-	frac35\log_2\tfrac35
=0.970951,\\
\entropy(\mathrm{Normal})&=-\tfrac45\log_2\tfrac45-	frac15\log_2\tfrac15
=0.721928,\\
\IG(Y,\mathrm{Humidity})
&=0.970951-\tfrac5{10}(0.970951)-\tfrac5{10}(0.721928)\\
&=\boxed{0.125}.
\end{align*}

\paragraph{Temperature.}
The Hot branch has $(1\ \mathrm{Yes},2\ \mathrm{No})$ and the Mild branch has
$(2\ \mathrm{Yes},1\ \mathrm{No})$, so both have entropy $0.918296$.  The Cool
branch has $(3\ \mathrm{Yes},1\ \mathrm{No})$, so its entropy is $0.811278$.
Therefore,
\begin{align*}
\IG(Y,\mathrm{Temperature})
&=0.970951-\tfrac3{10}(0.918296)-\tfrac3{10}(0.918296)
 -\tfrac4{10}(0.811278)\\
&=\boxed{0.095}.
\end{align*}
Since $0.322>0.125>0.095$, the first splitting attribute is
\(\boxed{\mathrm{Outlook}}\).

\subsection*{3(b) Force Humidity as the first split}

\subsubsection*{(i) Entropy in the two branches}
The branch entropies, already computed from their label counts, are
\[
\boxed{\entropy(Y\mid \mathrm{Humidity=High})=0.971},
\qquad
\boxed{\entropy(Y\mid \mathrm{Humidity=Normal})=0.722}.
\]

\subsubsection*{(ii) Information gain for Humidity}
The weighted post-split entropy is
\[
\tfrac5{10}(0.970951)+\tfrac5{10}(0.721928)=0.846439,
\]
so
\[
\boxed{\IG(Y,\mathrm{Humidity})=0.970951-0.846439=0.125}.
\]

\subsubsection*{(iii) Split each branch on Temperature}
For the High-humidity branch, Temperature=Hot contains $(1\ \mathrm{Yes},2\ \mathrm{No})$
and Temperature=Mild contains $(1\ \mathrm{Yes},1\ \mathrm{No})$.  Consequently,
\begin{align*}
\IG(\mathrm{High},\mathrm{Temperature})
&=0.970951-\tfrac35(0.918296)-\tfrac25(1.000000)\\
&=\boxed{0.020}.
\end{align*}

For the Normal-humidity branch, Temperature=Cool contains $(3\ \mathrm{Yes},1\ \mathrm{No})$
and Temperature=Mild contains one Yes label.  The latter branch is pure, so its
entropy is zero:
\begin{align*}
\IG(\mathrm{Normal},\mathrm{Temperature})
&=0.721928-\tfrac45(0.811278)-\tfrac15(0)\\
&=\boxed{0.073}.
\end{align*}

\subsection*{3(c) Two problems with decision trees}
\begin{enumerate}[label=\arabic*.]
    \item \textbf{Overfitting.} If the tree keeps splitting until leaves are
    nearly pure, it can memorize noise or rare training examples.  Its training
    error may be near zero while its test error is large.  Limiting depth,
    requiring a minimum leaf size, or pruning reduces this problem.
    \item \textbf{Instability under small data changes.} Greedy splitting makes
    an early choice that can change when one training point is added, removed,
    or slightly perturbed; the changed root then changes the entire subtree.
    Bagging or random forests reduce this variance by averaging many trees.
\end{enumerate}

\section{Regression}

The training points are $(7,6)$, $(2,2)$, and $(4,3)$.  With
$h_{\mathbf w}(x)=w_1x+w_0$, the sum-of-squared-errors objective is
\[
L(w_0,w_1)=(6-7w_1-w_0)^2+(2-2w_1-w_0)^2+(3-4w_1-w_0)^2.
\]

\subsection*{4(a) Least-squares weights}
Set the two partial derivatives to zero.  Equivalently, the normal equations are
\begin{align*}
3w_0+13w_1&=11,\\
13w_0+69w_1&=58,
\end{align*}
where
\[
\sum x_i=13,\quad \sum y_i=11,\quad
\sum x_i^2=49+4+16=69,\quad
\sum x_iy_i=42+4+12=58.
\]
The determinant is $3\cdot69-13^2=38>0$, so the quadratic objective has a
unique minimizer.  Eliminating $w_0$ gives
\[
(207-169)w_1=174-143
\quad\Longrightarrow\quad
38w_1=31
\quad\Longrightarrow\quad
w_1=\frac{31}{38}=0.815789.
\]
Substitution into the first normal equation gives
\[
3w_0=11-13\left(\frac{31}{38}\right)=\frac{15}{38},
\qquad
w_0=\frac{5}{38}=0.131579.
\]
Thus
\[
\boxed{h(x)=\frac{31}{38}x+\frac{5}{38}}.
\]

\subsection*{4(b) Predictions for all five points}
Using the same weights, the five predictions are
\begin{align*}
\widehat y_1=h(7)&=\frac{222}{38}=5.842105,\\
\widehat y_2=h(2)&=\frac{67}{38}=1.763158,\\
\widehat y_3=h(4)&=\frac{129}{38}=3.394737,\\
\widehat y_4=h(3)&=\frac{98}{38}=2.578947,\\
\widehat y_5=h(-1)&=-\frac{26}{38}=-0.684211.
\end{align*}
Therefore,
\[
\boxed{(\widehat y_1,\widehat y_2,\widehat y_3,\widehat y_4,\widehat y_5)
=(5.842105,1.763158,3.394737,2.578947,-0.684211)}.
\]

\subsection*{4(c) Overall loss on all five points}
The residuals $y_i-\widehat y_i$ are
\[
\frac{3}{19},\quad \frac{9}{38},\quad -\frac{15}{38},\quad
-\frac{11}{19},\quad -\frac{6}{19}.
\]
Therefore,
\begin{align*}
L_{1:5}
&=\left(\frac3{19}\right)^2+\left(\frac9{38}\right)^2
 +\left(\frac{15}{38}\right)^2+\left(\frac{11}{19}\right)^2
 +\left(\frac6{19}\right)^2\\
&=\frac{485}{722}=0.671745.
\end{align*}
Thus the requested overall loss is \(\boxed{0.671745}\).

\subsection*{4(d) $R^2$ statistic}
For all five observed outputs,
\[
\bar y=\frac{6+2+3+2-1}{5}=\frac{12}{5}=2.4,
\]
and the total sum of squares is
\begin{align*}
\mathrm{SST}&=\sum_{i=1}^{5}(y_i-\bar y)^2\\
&=(6-2.4)^2+(2-2.4)^2+(3-2.4)^2
 +(2-2.4)^2+(-1-2.4)^2\\
&=25.2.
\end{align*}
Using $\mathrm{SSE}=0.671745$ from part (c),
\begin{align*}
\Rtwo&=1-\frac{\mathrm{SSE}}{\mathrm{SST}}
=1-\frac{0.671745}{25.2}\\
&=\boxed{0.973343}\approx 97.3\%.
\end{align*}
On this five-point evaluation set, the model explains approximately $97.3\%$
of the variation in the outputs.  This is a strong fit for this small dataset;
it should not be interpreted as a guarantee of the same performance on a larger
or different population.

\end{document}
